\documentclass[11pt,a4paper]{article}
\usepackage[margin=27mm]{geometry}
\usepackage[T1]{fontenc}
\usepackage{lmodern}
\usepackage{amsmath,amssymb}
\usepackage{hyperref}
\hypersetup{colorlinks=true,linkcolor=blue,urlcolor=blue}
\setlength{\parindent}{0pt}
\setlength{\parskip}{7pt}
\title{Adams fixed points on a cyclic group\\Disproof of Conjecture 00000008230}
\author{ziangni-sys}
\date{4 October 2026}
\begin{document}
\maketitle

\textbf{Result.} For the cyclic group $C_3$ and the second Adams operation,
the fixed subspace of complex class functions has dimension $2$, whereas
the number of $2$-regular conjugacy classes is $3$. Thus the dimension
equality asserted in the conjecture is false.

\section*{1. Definitions and the disputed assertion}
For a finite group $G$, write $\mathrm{CF}(G;\mathbb C)$ for its vector
space of complex class functions. These are functions $f:G\to\mathbb C$
such that $f(xgx^{-1})=f(g)$ for every $x,g\in G$.
The Adams operation on class functions is
\[
  (\psi^k f)(g)=f(g^k).
\]
Its fixed subspace consists of the class functions satisfying
$f(g^k)=f(g)$ for all $g$. An element is $k$-regular when its order is
coprime to $k$; a conjugacy class is $k$-regular when its elements are.
For $k=2$, this is simply the familiar odd-order condition.

The English statement explicitly asserts that the dimension of this
fixed subspace equals the number of $k$-regular classes. The Chinese
statement asserts the same dimension equality. We disprove this conjunct.
The conjecture's other assertions about spectral support and composition
corrections are not needed for the contradiction.

\section*{2. The counterexample}
Let
\[
  G=C_3=\langle a\mid a^3=1\rangle=\{1,a,a^2\},\qquad k=2.
\]
The group is abelian, so its three conjugacy classes are the singletons
$\{1\}$, $\{a\}$, and $\{a^2\}$. Their element orders are $1$, $3$,
and $3$, respectively. All three are coprime to $2$. Consequently,
\[
  \#\{\text{$2$-regular conjugacy classes of $C_3$}\}=3.
\]
Every function on $C_3$ is a class function. Writing its values in the
order $(1,a,a^2)$, the second Adams operation is
\[
  \psi^2(u,v,w)=(u,w,v),
\]
since $1^2=1$, $a^2=a^2$, and $(a^2)^2=a^4=a$.

\newpage
Thus $\psi^2f=f$ holds exactly when $f(a)=f(a^2)$. Its fixed subspace is
\[
  F=\{(u,v,v):u,v\in\mathbb C\}.
\]
The maps
\[
  F\longrightarrow\mathbb C^2,\quad f\longmapsto(f(1),f(a)),
  \qquad
  (u,v)\longmapsto\bigl[g\mapsto
    \begin{cases}u,&g=1,\\v,&g\ne1\end{cases}\bigr]
\]
are inverse complex linear maps. Therefore $\dim_{\mathbb C}F=2$.
Equivalently, the functions with values $(1,0,0)$ and $(0,1,1)$ form a
basis. Combining the two exact computations gives
\[
  \boxed{\dim_{\mathbb C}\ker(\psi^2-I)=2\ne3
    =\#\{\text{$2$-regular conjugacy classes}\}.}
\]
This disproves the stated equality and hence the conjunctive conjecture.
\hfill$\square$

The reason the proposed count fails is already visible here: squaring
permutes the two nonidentity conjugacy classes, so a fixed function must
take equal values on them. Counting the classes individually does not
count the independent values of a fixed function.

\section*{3. Lean correspondence and verification}
The Lean project uses Lean 4.19.0 and pins Mathlib to commit
\begin{center}\small
\texttt{c44e0c8ee63ca166450922a373c7409c5d26b00b}.
\end{center}
The finite group is \texttt{Multiplicative (ZMod 3)}, with its actual
group operations. Conjugacy classes are Mathlib's quotient
\texttt{ConjClasses G}; regularity uses \texttt{orderOf} and
\texttt{Nat.Coprime}. In particular, neither the number of classes nor
the element orders are supplied as unproved data.

\texttt{cube\_one} checks the group identity $g^3=1$ for every element.
It follows that each order divides $3$, proving regularity.
\texttt{regular\_class\_count} computes the cardinality of the actual
subtype of regular conjugacy classes, using the conjugacy-class
bijection for a commutative group.

\texttt{fixedClassFunctions} is a complex submodule whose defining
conditions are conjugation invariance and the equation $\psi^2f=f$.
The linear equivalence \texttt{coordinates} implements the two inverse
maps above and proves both inverse identities and linearity.
\texttt{fixed\_dimension} then obtains dimension $2$ through this
equivalence, rather than assuming a basis or a dimension table.

The final theorem \texttt{Adams8230.counterexample} states that the two
actual quantities are unequal. Reproduction commands and the axiom audit
are in \texttt{VERIFICATION.md}. There are no additional axioms,
incomplete proofs, or \texttt{native\_decide}; no auxiliary computational
program is required.
\end{document}
